\guidelinesubsubsection{Rationale}

Even with well-justified benchmarks and metrics (\benchmarksmetrics), automated measurement captures only the constructs the benchmark was designed to measure.
Subjective constructs and constructs that no current benchmark covers require human assessment.

\guidelinesubsubsection{Recommendations}

When LLMs automate research or software development tasks previously performed by humans, the LLM's performance needs to be assessed.
Where no benchmark adequately operationalizes the target construct, researchers \should validate LLM outputs against human judgment.

\guidelineparagraph{Study Design Considerations}
Studies that include human participants need additional considerations, including a recruitment strategy, annotation guidelines, training sessions, or ethical approvals.
Therefore, researchers \should consider human validation early in the study design, not as an afterthought.
Authors \must clearly define in the \paper the constructs that the human and LLM annotators evaluate~\cite{DBLP:conf/ease/RalphT18}.
% GROUNDING DBLP:conf/ease/RalphT18: "many software engineering studies simply assume that proposed measures are valid and make no attempt to assess construct validity"
When designing custom instruments to assess LLM output (e.g., questionnaires, scales), researchers \must share these instruments.

\guidelineparagraph{Replacing Human Judgment}
LLMs may be used to replace, rather than augment, humans in research tasks such as annotating software artifacts, coding interview transcripts, or simulating participants in user studies. In such cases, researchers \must explain whether and how the replacement is justified and \should follow systematic approaches to support this judgment, e.g., showing that a jury of three LLMs exhibits inter-model agreement at the same threshold researchers would require for human-to-human agreement, such as Krippendorff's $\alpha>0.8$. Such agreement is necessary but not sufficient. High inter-model agreement can reflect shared model biases rather than valid annotation, so researchers \should additionally validate a sample against human experts before treating LLM outputs as a substitute for human work~\cite{DBLP:conf/msr/AhmedDTP25, Krippendorff2018}. For qualitative coding of interpretive data such as open-ended developer interview responses, agreement metrics do not show that an LLM can perform the interpretive work that the task requires. Researchers \should justify why the LLM is methodologically appropriate (see \llmsforresearcher).
% GROUNDING DBLP:conf/msr/AhmedDTP25: "replacing more than one human can inflate inter-rater agreement because model-model agreements are much higher"
% GROUNDING Krippendorff2018: "Reliability does not guarantee validity."

\guidelineparagraph{Subjective Judgment and Agreement}
When the judgment is subjective (i.e., depends on the judge's values or theories), the same artifacts \should be judged independently by both the LLM and a panel of human experts, and the LLM judgments \should then be compared against an aggregated human reference. Researchers \should clearly describe their aggregation method and reasoning.
Researchers \should use established reference models to compare humans with LLMs.
For example, \citeauthor{DBLP:conf/icse/SchneiderFW25} outline design considerations for studies comparing LLMs with humans~\cite{DBLP:conf/icse/SchneiderFW25}.
% GROUNDING DBLP:conf/icse/SchneiderFW25: "We propose a framework to support researchers in designing and executing studies, and in assessing LLMs with respect to humans."

One systematic approach to building the human reference is to randomly order the objects requiring judgment and put them in groups small enough to judge in one to three hours. Two or three human experts review the first group together, simultaneously judging, discussing, and creating a set of decision rules documenting their reasoning. Then, the experts iterate among rounds of:
\begin{enumerate}
    \item independently rating one group of objects,
    \item calculating inter-rater agreement (IRA) or reliability (IRR),
    \item meeting to discuss disagreements and reach consensus,
    \item updating the decision rules so the disagreement will not recur.
\end{enumerate}

The goal is to reach a target IRA or IRR (e.g., Krippendorff's $\alpha>0.8$), indicating sufficient decision rules. Once this target is reached and sustained for two to three groups, it is permissible to continue with a single rater.
Researchers \should report measures of IRA or IRR, preferably broken down by round, and \should apply thresholds consistent with established standards. \citeauthor{Krippendorff2018} recommends discarding data with $\alpha<0.667$, considers $0.667 \leq \alpha < 0.8$ sufficient only for tentative conclusions, and requires $\alpha \geq 0.8$ for reliable data~\cite{Krippendorff2018}.
% GROUNDING Krippendorff2018: "Rely only on variables with reliabilities α ≥ 0.800. Consider variables with reliabilities between α = 0.667 and α = 0.800 only for drawing tentative conclusions."

Confounding factors \should be discussed and, where feasible, controlled for (e.g., by categorizing participants according to their level of experience or expertise).
For value-laden or culturally contingent constructs (e.g., judging code-style appropriateness or comment helpfulness), researchers \should describe rater demographics beyond expertise (e.g., geographic, linguistic, or professional background) and discuss potential demographic biases in rater recruitment and instructions~\cite{bean2025measuring}.
% GROUNDING bean2025measuring: "If using human raters, describe their demographics and mitigate potential demographic biases in rater recruitment and instructions"
Where applicable, researchers may perform a power analysis~\cite{Cohen1992,DBLP:journals/infsof/DybaKS06} to estimate the required sample size, ensuring sufficient statistical power in their experimental design.
% GROUNDING Cohen1992: "The sample sizes necessary for .80 power to detect effects at these levels are tabled for 8 standard statistical tests"
% GROUNDING DBLP:journals/infsof/DybaKS06: "Only one of the papers that performed an a priori power analysis used it to guide the choice of sample size."
Although established sample-size guidance for LLM-human comparisons is limited, related fields commonly use 100 comparisons without further justification~\cite{DBLP:journals/npjdm/TamSKSPMOWVFMCSPW24}.
% GROUNDING DBLP:journals/npjdm/TamSKSPMOWVFMCSPW24: "The majority of studies have 100 or below LLMs output for human evaluation"

\guidelineparagraph{Agentic Tools}
Agentic human-in-the-loop interaction is a built-in human validation, with larger degrees of freedom than traditional experiments that use LLMs directly.
Besides generating and modifying content, agentic systems such as \emph{Claude Code} (see \design) can autonomously call command-line tools or pull in additional information from MCP servers.
When evaluating agentic tools, researchers \should assess the feedback that users provided on the agent's proposed actions (e.g., file edits, command executions), report statistics on how frequently they accepted the proposals, and how they modified them.

\guidelinesubsubsection{Examples}

\citeauthor{DBLP:conf/msr/AhmedDTP25} proposed a systematic method for deciding whether LLMs can replace human annotators on a given task, using inter-model agreement as an initial screening criterion (with a threshold of $\alpha>0.5$) and model confidence for sample-level decisions~\cite{DBLP:conf/msr/AhmedDTP25}. However, their threshold is well below the levels generally considered acceptable for inter-rater agreement.
% GROUNDING DBLP:conf/msr/AhmedDTP25: "model confidence as a means to select specific samples where LLMs can safely replace human annotators"
Notably, while three LLMs exhibited higher inter-model agreement than human annotators on some tasks, human-model agreement remained low on others.
This illustrates that high inter-model reliability does not guarantee alignment with human judgment, reinforcing the need for human validation before assuming that LLM annotations are valid.
Moreover, a high inter-model agreement could merely indicate that the models share systematic biases and hence reliably agree on the wrong answer.

\citeauthor{DBLP:journals/corr/abs-2501-19297} evaluated ChatGPT's capability to generate requirements documents by comparing an LLM-generated and a human-generated document based on the same business use case~\cite{DBLP:journals/corr/abs-2501-19297}. Domain experts reviewed both documents and attempted to distinguish their origin.
% GROUNDING DBLP:journals/corr/abs-2501-19297: "Document "01": Created by an LLM, ChatGPT-4.0 using a single-shot prompt. Document "02": Created by a human expert during a one-hour, closed session."
For agentic tools, \citeauthor{DBLP:conf/icse-seip/TakerngsaksiriPTTZJLCCW25} reported acceptance and modification rates from real users for \emph{HULA}, an agentic plan-and-code system deployed in Atlassian JIRA~\cite{DBLP:conf/icse-seip/TakerngsaksiriPTTZJLCCW25}.
% GROUNDING DBLP:conf/icse-seip/TakerngsaksiriPTTZJLCCW25: "practitioners approved the generated plans for 433 out of 527 plan generated issues, achieving a plan approval rate of 82%"

\guidelinesubsubsection{Benefits}

Validation against human judgments builds confidence in the LLM's accuracy and validity~\cite{khraisha2024canlargelanguagemodelshumans}. Reporting IRA or IRR supports this validation but does not stand in for it, since even high agreement may reflect shared biases rather than valid judgment. When human reviewers disagree with the LLM, the disagreement points to concrete improvements in the prompts, the context provided to the LLM, or the operational definition of the construct.
% GROUNDING khraisha2024canlargelanguagemodelshumans: "TP and TN are the cases where GPT-4 agreed with human reviewers, meaning it made correct decisions."

\guidelinesubsubsection{Challenges}

Assessing LLM performance with respect to a construct such as code maintainability, answer correctness, or annotation reliability is challenging because ensuring that a construct is defined well and operationalized using an appropriate measurement model requires a deep understanding of (1) the construct, (2) construct validity in general, and (3) instrumentation~\cite{DBLP:journals/tse/SjobergB23,DBLP:books/sp/24/RalphKAA24}. Comparing an LLM to human judges is typically slower and more expensive than machine-generated measures. More fundamentally, neither human judgment nor machine-generated measures provides an objective ground truth against which LLM accuracy can be firmly determined. Human preference ratings under-represent properties such as factuality and faithfulness, and assertively phrased outputs receive systematically fewer flagged errors from crowdworkers~\cite{DBLP:conf/iclr/HoskingBB24}. Human validation is worth its added cost only when automated measures alone fail to capture the constructs the study evaluates.
% GROUNDING DBLP:journals/tse/SjobergB23: "CV requires an unambiguous definition or an established understanding of the concept involved in the construct. Without such a basis, it is difficult to evaluate how well a set of indicators represents the concept."
% GROUNDING DBLP:books/sp/24/RalphKAA24: "You need to think deeply about what exactly you're trying to measure. Ask yourself not only 'what am I trying to measure?' but also 'from whose perspective?' and 'in what context?'."
% GROUNDING DBLP:conf/iclr/HoskingBB24: "annotators are more trusting of assertive responses, and are less likely to identify factuality or inconsistency errors within them"

Human judgments exhibit variability due to differences in experience, expertise, interpretations, and personal biases~\cite{DBLP:journals/pacmhci/McDonaldSF19}. In pairwise judgment of LLM outputs, both human and LLM judges shift their preferences toward answers carrying fake references or richly formatted content, regardless of correctness~\cite{DBLP:conf/emnlp/ChenCLJW24}. When diverse humans rate items reliably given clear decision rules, we assume that reliability implies validity, but it does not. Measuring reliability is \textit{much} easier than measuring validity, and often the best researchers can do is argue conceptually for why their judges, decision rules, and constructs \textit{should} produce valid ratings.
% GROUNDING DBLP:journals/pacmhci/McDonaldSF19: "they may be interpreted differently by different people depending on researchers' backgrounds and the focus of the analysis"
% GROUNDING DBLP:conf/emnlp/ChenCLJW24: "it refers to the inclination that judges tend to prefer visually appealing content, regardless of its actual validity."

\guidelinesubsubsection{Study Types}

This guideline applies to all study types, although the need for human validation varies. \emph{Replacing Human Judgment} introduces a conditional \must that applies particularly to \annotators, \judges, \synthesis, and \subjects, where the LLM substitutes for human input.
For \annotators, researchers \should validate LLM-generated annotations against human annotators to assess labeling quality and identify systematic biases.
When using \judges, researchers \should co-create initial rating criteria with humans and validate a sample of LLM judgments against human expert assessments.
For \synthesis, researchers \should employ human oversight to verify that qualitative interpretations and synthesized outputs faithfully represent the underlying data.
For \subjects, researchers \should validate simulated responses against real human data to assess the fidelity of the simulation.
For \llmusage, researchers \should carefully reflect on the validity of their evaluation criteria and validate subjective assessments with human experts.
For \newtools whose output is supposed to match human expectations, researchers \should validate the LLM output against human judgment.
For \benchmarkingtasks, there is less need for human validation when using extensively validated and widely-used benchmarks, but researchers \should employ human validation when creating or adapting new benchmarks.

\guidelinesubsubsection{Advice for Reviewers}

Human validation may be the most challenging of our guidelines to assess because it often requires evaluating conceptual arguments. If LLM output is validated only by comparison with other LLMs, reviewers should look for \textit{quantitative empirical evidence} that such comparison is reliable \textit{and} valid. High inter-model agreement alone is insufficient, as reliability does not imply validity. Similarly, reviewers should expect evidence that any employed benchmarks are reliable and valid. Absent such evidence, human validation is warranted.
A single human judge is appropriate only when judgments depend on widely accepted theories and involve limited value conflict (e.g., tagging method names containing abbreviations). For multiple judges, reviewers should expect IRA/IRR improvement techniques as described in the recommendations above (experienced raters, organized rounds, consensus meetings, updated decision rules). Low IRA or IRR (e.g., Krippendorff's $\alpha<0.8$) without these techniques is a concern. Conversely, if authors have followed best practices and still obtained mediocre results (e.g., $0.66<\alpha<0.8$), this should be noted as a limitation.
Beyond reliability, reviewers should expect authors to explain conceptually why their human judgments should be valid, considering construct definitions, decision rules, and judge expertise.
As with other guidelines, missing information is typically a revision request. Absent judge instructions, instruments, decision rules, or construct definitions may prevent assessment of rigor and validity, while missing recruitment details are less critical. Clarification requests about construct definitions are routine and should not alone warrant rejection.

\guidelinesubsubsection{See Also}
\begin{itemize}[label=$\rightarrow$]
    \item \seebenchmarksmetrics: Human evaluation complements automated metrics where benchmarks cannot sufficiently capture the target construct.
    \item \seemodelversion: Human-validation findings characterize only the specific model and configuration that produced the outputs.
    \item \seedesign: Prompt and architecture choices produce the outputs that humans then validate.
    \item \seetraces: Stored interaction logs and runtime traces are the artifacts human reviewers examine.
    \item \seelimitationsmitigations: Rater demographics, threshold choices, and the validity of human judgments are limitations to discuss.
\end{itemize}